\begin{exercisechunk}
\item \textbf{The cost of polynomial beta-mixing.}
\label[exercise]{ex:l6-polynomial-mixing}
\exerciseprofile{2}{2}{2}{\exproof\quad\excalculation\quad\exinterpretation}{%
Convert a polynomial mixing rate into a sample-size bound and compare its
dependence on accuracy with iid and geometric-mixing bounds.}
Assume the setting of \cref{thm:mixing-finite-class} and
\(\beta(a)\leq Ca^{-r}\) for constants \(C,r>0\). Let
\(\epsilon,\delta\in(0,1)\) and \(A=\log(2|\cF|/\delta)\).
\begin{enumerate}[(a),beginpenalty=10000]
\item Starting from \cref{eq:mixing-bound}, give sufficient conditions on
the selected-sample size \(q\) and gap \(a\) that make each term at most
\(\delta/2\).
\item Use \(n\geq1+(q-1)a\) and a suitable choice of \(q,a\) to show, up to
universal constants and integer rounding, that it suffices to take
\[
 n\gtrsim \frac{A}{\epsilon}
 +\left(\frac{A}{\epsilon}\right)^{1+1/r}
  \left(\frac{C}{\delta}\right)^{1/r}.
\]
\item Compare the dependence on \(\epsilon\) with iid sampling and with
geometric beta-mixing. Is the degradation exponential in \(1/\epsilon\)?
\end{enumerate}
\begin{hints}
Use \((1-\epsilon)^q\leq e^{-\epsilon q}\). First choose
\(q\asymp A/\epsilon\), then choose \(a\) so that
\(qCa^{-r}\lesssim\delta\), and finally take \(n\) of order \(qa\).
\end{hints}
\end{exercisechunk}
